\section{Naturalism and Interwar Germany}

After the devastations of World War II, Evola --- who was never a Fascist nor a National Socialist --- turned to an analysis of those two movements from the perspective of Tradition in his publication
\emph{Fascism viewed from the right with notes on the III Reich}.


Both movements suffered from a confusion of Traditional elements with ideals from the Enlightenment. Insofar as they were dominated by materialism, economic concerns, and scientism, they lacked a transcendent --- and hence, Traditional --- perspective.

As Evola pointed out in ``The Individual and the Becoming of the World''\footnote{\url{https://www.amazon.com/dp/0999086049}}, the magical act consists in bringing the idea (essence) into manifestation (existence). If the idea is confused or incoherent, the result will be disorder. In the following passage, Evola points out the strong streak of naturalism that dominated the ideology of National Socialism. As such, even if unconsciously held, it must eventually serve the aims and purposes of organized naturalism.

For Evola, everything is ultimately a spiritual problem, impervious to the methods of secular science. All depends on the quality of men, not on the layout of their chromosomes.

\begin{quotex}
As for Hitler himself, valid contributions to the problematic of the vision of the world in a higher sense are not found in his writing and speeches. His Wagnerian infatuation is already significant --- for him, as for Chamberlain, Richard Wagner is valued as ``the prophet of Germanism'' --- his incapacity to recognize the measure in which, apart from the greatness of his romantic art, the distortions of not a few German and Nordic traditions and sagas must be attributed to Wagner. If in his writings and speeches Hitler often mentions God and Providence, of which he considered himself the designated one and executor, it is not clear what this Providence could be if he, on the one hand, following Darwin somewhat more than Nietzsche, recognized the right of the strongest as the supreme law of life, and on the other excluded as superstition any intervention or supernatural order whatsoever and gave himself over to an exaltation of modern science and the ``eternal laws of nature''. Such an attitude was likewise characteristic of the principle ideologue of the movement,Alfred Rosenberg, who came to see in modern science ``a creation of the pure Aryan spirit'', without taking into consideration that if technical conquests are owed to it, one also owes to them the more deleterious and irreversible spiritual devastations of the modern era, the desacralisation of the universe. And a bluntly Enlightenment and rationalist incomprehension of the essential aspects of religion, paradoxically going hand in hand with the mysticism of blood, if it appeared in Hitler, it was fully explicit in Rosenberg; rites and sacraments for him were superstition, creations of a non-Aryan spirit.
\flright{\textsc{Julius Evola}, \textit{Notes on the Third Reich}}
\end{quotex}


\flrightit{Posted on 2008-08-24 by Cologero}

